\subsection{SUBRINGS AND IDEALS IN A QUOTIENT RING}

PROPOSITION 5. Let A and A' be two rings, f a homomorphism of A into A' and a the kernel of f.

\begin{enumerate}
\def\labelenumi{(\alph{enumi})}
\item
  Let \(\mathbf{B}'\) be a subring of \(\mathbf{A}'\) . Then \(\mathbf{B} = f(\mathbf{B}')\) is a subring of \(\mathbf{A}\) containing \(\mathfrak{a}\) . If \(f\) is surjective, then \(f(\mathbf{B}) = \mathbf{B}'\) and \(f \mid \mathbf{B}\) defines when passing to the quotient an isomorphism of \(\mathbf{B}/\mathfrak{a}\) onto \(\mathbf{B}'\) .
\item
  Let \(\mathfrak{b}'\) be a left (resp. right, two-sided) ideal of \(\mathbf{A}'\) . Then \(\mathfrak{b} = f(\mathfrak{b}')\) is a left (resp. right, two-sided) ideal of \(\mathbf{A}\) containing \(\mathfrak{a}\) .
\item
  If \(\mathfrak{b}'\) is a two-sided ideal of \(\mathbf{A}'\) , the composite mapping of the canonical morphism \(\mathbf{A}' \to \mathbf{A}' / \mathfrak{b}'\) and \(f: \mathbf{A} \to \mathbf{A}'\) defines, when passing to the quotient, an injective morphism \(\bar{f}\) of \(\mathbf{A} / \mathfrak{b}\) into \(\mathbf{A}' / \mathfrak{b}'\) . If \(f\) is surjective, \(\bar{f}\) is an isomorphism of \(\mathbf{A} / \mathfrak{b}\) onto \(\mathbf{A}' / \mathfrak{b}'\) .
\item
  Suppose \(f\) is surjective. Let \(\Phi\) be the set of subrings (resp. left ideals, right ideals, two-sided ideals) of A containing a. Let \(\Phi'\) be the set of subrings (resp. left ideals, right ideals, two-sided ideals) of \(A'\) . The mappings \(B \to f(B)\) and \(B' \mapsto f'(B')\) are inverse bijections of \(\Phi\) onto \(\Phi'\) and \(\Phi'\) onto \(\Phi\) .
\end{enumerate}

(a) and (b) are obvious, except the last assertion of (a) which follows from no. 7, Theorem 3.

The composite morphism \(g: \mathbf{A} \to \mathbf{A}' \to \mathbf{A}' / \mathfrak{b}'\) considered in (c) has kernel \(\mathfrak{b}\) and hence \(\bar{f}\) is an injective morphism of \(\mathbf{A} / \mathfrak{b}\) into \(\mathbf{A}' / \mathfrak{b}'\) (§ 8, no. 7, Theorem 3). If \(f\) is surjective, \(g\) is surjective and hence \(\bar{f}\) is surjective.

Suppose \(f\) is surjective. By the above, the mapping \(\theta: \mathbf{B}' \mapsto f(\mathbf{B}')\) is a mapping of \(\Phi'\) into \(\Phi\) . Clearly the mapping \(\eta: \mathbf{B} \mapsto f(\mathbf{B})\) is a mapping of \(\Phi\) into \(\Phi'\) . Then \(\theta \circ \eta = \operatorname{Id}_{\Phi}, \eta \circ \theta = \operatorname{Id}_{\Phi'}\) , whence (d).

Remark. In the above notation, \(\theta\) and \(\eta\) are ordered set isomorphisms ( \(\Phi\) and \(\Phi'\) being ordered by inclusion).

COROLLARY. Let A be a ring and a a two-sided ideal of A.

\begin{enumerate}
\def\labelenumi{(\alph{enumi})}
\item
  Every left (resp. right, two-sided) ideal of A/a can be written uniquely in the form b/a, where b is a left (resp. right, two-sided) ideal of A containing a.
\item
  If \(\mathfrak{b}\) is two-sided, the composite homomorphism \(\mathbf{A} \to \mathbf{A} / \mathfrak{a} \to (\mathbf{A} / \mathfrak{a}) / (\mathfrak{b} / \mathfrak{a})\) defines when passing to the quotient an isomorphism of \(\mathbf{A} / \mathfrak{b}\) onto \((\mathbf{A} / \mathfrak{a}) / (\mathfrak{b} / \mathfrak{a})\) .
\end{enumerate}

It suffices to apply Proposition 5 to the canonical morphism of A onto A/a.

